\scanpage{28}
% archive sheet number: 27
$P^1_S$ sections $0,1,\infty$. Le groupe des automorphismes s'identifie au groupe $\unsure{\mathfrak{S}|\mathfrak{G}}_3$ des permutations des trois objets $0,1,\infty$, groupe \unsure{à|:} 6 éléments $1,\varphi,\psi,\lambda,\mu,\nu$

\begingroup\makeatletter\tagsleft@true\makeatother
\[
\left\{\begin{aligned}
\varphi(t)&=1-\frac1t\\
\psi(t)&=\frac1{1-t}
\end{aligned}\right.\qquad
\text{2 permutations circulaires non triviales sur }0,1,\infty
\tag*{(1)}
\]
\endgroup

\begingroup\makeatletter\tagsleft@true\makeatother
\[
\left\{\begin{array}{ll||l}
\lambda t=\frac1t&\text{transposition de }0,\infty&1\\
\mu t=1-t&\text{transposition de }0,1&\infty\\
\nu t=1+\frac1{t-1}=\frac{t}{t-1}&\text{transposition de }1,\infty&0
\end{array}\right.
\tag*{(2)}
\]
\endgroup
% slanted
\margin{\unsure{renommer|renumérer|prenons|pour mémoire|pour moi} :\par $\left\{\begin{aligned}\lambda t&=1-\unsure{t|1}\\ \mu t&=\tfrac{t}{t-1}\\ \nu t&=1/t\end{aligned}\right.$\par ??}

\begingroup\makeatletter\tagsleft@true\makeatother
\[
\boxed{\begin{aligned}\varphi^2&=\psi\qquad\psi^2=\varphi\\
\deleted{\unsure{\psi\varphi|\varphi\psi}}\quad\varphi\psi&=\psi\varphi=1\end{aligned}}
\tag*{3 a)}
\]
\endgroup

\begingroup\makeatletter\tagsleft@true\makeatother
\[
\boxed{\lambda^2=\mu^2=\nu^2=1}\tag*{b)}
\]
\endgroup

\begingroup\makeatletter\tagsleft@true\makeatother
\[
\boxed{\begin{aligned}\lambda\mu&=\mu\nu=\nu\lambda=\varphi\\
\mu\lambda&=\nu\mu=\lambda\nu=\psi\end{aligned}}\tag*{c)}
\]
\endgroup

\begingroup\makeatletter\tagsleft@true\makeatother
\[
\boxed{\begin{array}{l|l}
\varphi\lambda=\mu\varphi=\nu&\psi\nu=\mu\psi=\lambda\\
\varphi\mu=\nu\varphi=\lambda&\psi\mu=\lambda\psi=\nu\\
\varphi\nu=\lambda\varphi=\mu&\psi\lambda=\nu\psi=\mu
\end{array}}\tag*{d)}
\]
\endgroup

$1,\varphi,\psi$ groupe cyclique d'ordre 3, de \unsure{générateurs|engendré par} $\varphi$ et $\psi$, invariant dans $\mathfrak{S}_3$

Le quotient est $\mathbb{Z}/2$ et $\mathfrak{S}_3$ est produit semi-direct ; il y a trois sous-groupes \unsure{sections|section}, engendrés resp\textsuperscript{t}. par $\lambda,\mu,\nu$ (les seuls \unsure{éléments} d'ordre 2)

\begingroup\makeatletter\tagsleft@true\makeatother
\[
\left\{\begin{aligned}
{}^\varphi\lambda&=\mu&{}^\varphi\mu&=\nu.&{}^\varphi\nu&=\lambda\\
{}^\psi\nu&=\mu&{}^\psi\mu&=\lambda&{}^{\psi}\lambda&=\nu
\end{aligned}\right.\tag*{4)}
\]
\endgroup

\begingroup\makeatletter\tagsleft@true\makeatother
\[
\left\{\begin{aligned}
{}^\lambda\varphi&={}^\mu\varphi={}^\nu\varphi=\psi\\
{}^\lambda\psi&={}^\mu\psi={}^\nu\psi=\varphi
\end{aligned}\right.\tag*{5)}
\]
\endgroup

De plus, on a dans $P^1$ un automorphisme $\unsure{\sigma|5}$ \deleted{\unsure{de|dans}} de carré 1, donné par
\begingroup\makeatletter\tagsleft@true\makeatother
\[
\sigma t=\frac{t-2}{2t-1}\qquad[\sigma^2=\mathrm{id}]\tag*{6)}
\]
\endgroup
\unsure{permutant :|permutant à|permutable à|Normalisateur :} $\mathfrak{S}_{6}$ (c'est d'ailleurs le seul, sauf

